% Л10 -- DI, исходящие вызовы и конфигурация. Slides.
%
% Companion to the spoken script in ../L10-di-and-configuration.md, which is
% the source of truth: the slides carry diagrams, code, real terminal output
% and short marks, and the delivery lives in the script only -- this deck
% carries no speaker notes. Slide text is in Russian because that is the
% content; everything else in this file is English.
%
% Build:  make          (latexmk -pdf; engine: pdflatex)
%
% Every C# listing was compiled in a scratch `dotnet new web -n HelloWorld`;
% terminal output on the slides is captured from runs of that project. See
% ../CLAUDE.md for the verification table.
%
% House rules: ../CLAUDE.md and ../../../CLAUDE.md. Theme: ../../../../theme/.

\documentclass[aspectratio=169,10pt,t]{beamer}

\usetheme{dotnet}
\usetikzlibrary{arrows.meta,positioning,calc,decorations.pathreplacing,fit}
\usepackage{array}   % ragged-right table columns

% A ragged-right paragraph column: justified cells in narrow columns hyphenate
% badly, and a split word is the one thing the back row cannot read.
\newcolumntype{L}[1]{>{\raggedright\arraybackslash}p{#1}}

\dotnetlecture{Л10}

\title[Л10. DI, исходящие вызовы и конфигурация]{DI, исходящие вызовы и конфигурация}
\author{Дмитрий Курочкин}
\date{}

% Code sizes, one step down from \footnotesize. Smaller than 6.6pt does not
% carry to the back of the room.
\newcommand{\codesm}{\ttfamily\fontsize{7.6}{9.0}\selectfont}
\newcommand{\codexs}{\ttfamily\fontsize{7.0}{8.3}\selectfont}
\newcommand{\codexxs}{\ttfamily\fontsize{6.8}{7.8}\selectfont}
\newcommand{\termsm}{\ttfamily\fontsize{6.6}{7.9}\selectfont}

% A mark at the start of a listing line, for labels set beside the listing.
% Used through the listing's escapechar; the labels are drawn by an overlay
% picture after the listing, so the build needs two passes (latexmk does that).
\newcommand{\mk}[1]{\tikz[remember picture,overlay]\coordinate (#1) at (0,0.55ex);}

\tikzset{
  rawbrace/.style={decorate,decoration={brace,amplitude=4pt},hairline,
                   line width=0.8pt},
  rawlab/.style={anchor=west,font=\footnotesize,inner sep=0pt},
  box/.style={draw=hairline,line width=0.8pt,fill=codebg,inner sep=4pt,
              align=center,minimum height=0.9cm},
  arr/.style={-{Straight Barb[length=1.6mm]},muted,line width=0.7pt},
  arra/.style={-{Straight Barb[length=1.8mm]},accent,line width=0.9pt},
  arrl/.style={-{Straight Barb[length=1.8mm]},alt,line width=0.9pt},
}

\begin{document}
\shorthandoff{"}

% ------------------------------------------------------------------ title --
% 1
\begin{frame}[plain]
  \titlepagednc{Лекция 10}{DI, исходящие вызовы и конфигурация}{}{Дмитрий Курочкин}
\end{frame}

% ================================================ 1. DEPENDENCY INJECTION ==
\section{1. Внедрение зависимостей}

% 2
\begin{frame}[plain]
  \sectionopener{1}{Внедрение зависимостей}{Сервис не знает, \hi{из чего} он
    состоит.}
\end{frame}

% 3
\begin{frame}[plain]
  \vote{Сервису заказов нужно хранилище. Что вы пишете в конструкторе?}
\end{frame}

% 4
\begin{frame}[fragile]{Было}
  \slidetop
\begin{lstlisting}[style=csharp,basicstyle=\codesm,escapechar=|]
public sealed class OrderServiceOld
{
    private readonly SqlOrderRepository _repository =
|\mk{o1}|        new("Server=prod;Database=orders");
    public Task<Order?> FindAsync(int id, CancellationToken ct) =>
        _repository.FindAsync(id, ct);
}
\end{lstlisting}
  \begin{tikzpicture}[remember picture,overlay,every node/.style={rawlab}]
    \node[text=alt] at ([xshift=8.6cm]o1) {«я знаю, из чего состою»};
  \end{tikzpicture}

  \rulegap
  \begin{cbul}{alt}
    \item сервис \lo{прибит} к SQL-хранилищу
    \item строка подключения в коде: другая база --- пересборка
    \item тест требует настоящую базу
    \item десять сервисов --- десять \texttt{new}, десять экземпляров
  \end{cbul}
\end{frame}

% 5
\begin{frame}{Три слова}
  \slidetop
  \claim{Принцип можно соблюдать и \hi{без контейнера}.}

  {\renewcommand{\arraystretch}{1.6}
  \begin{tabular}{@{}L{4.6cm}L{8.2cm}@{}}
    \hi{\textbf{инверсия управления}}    & принцип: класс объявляет потребность,
      удовлетворяет её кто-то снаружи \\
    \hi{\textbf{внедрение зависимостей}} & способ: зависимость приходит через
      конструктор \\
    \hi{\textbf{контейнер}}              & инструмент: собирает объекты за вас \\
  \end{tabular}}

  \slidegap
  \aside{Пять классов собирают руками в \texttt{Main}. Контейнер нужен, когда
    классов сто.}
\end{frame}

% 6
\begin{frame}[fragile]{Стало}
  \slidetop
\begin{lstlisting}[style=csharp,basicstyle=\codesm,escapechar=|]
public sealed class OrderService(
|\mk{n1}|    IOrderRepository repository)
{
    public Task<Order?> FindAsync(int id, CancellationToken ct) =>
        repository.FindAsync(id, ct);
}

|\mk{n2}|// composition root
builder.Services.AddScoped<IOrderRepository, SqlOrderRepository>();
builder.Services.AddScoped<OrderService>();
\end{lstlisting}
  \begin{tikzpicture}[remember picture,overlay,every node/.style={rawlab}]
    \node[text=accent] at ([xshift=8.6cm]n1) {говорит, что ему нужно};
    \node at ([xshift=8.6cm]n2) {единственное место, где виден состав};
  \end{tikzpicture}

  \slidegap
  \aside{Кто создаст реализацию, какую и с какой строкой подключения ---
    сервис \hi{не знает}. Состав описан до \texttt{Build}.}
\end{frame}

% 7
\begin{frame}[fragile]{Формы регистрации}
  \slidetop
\begin{lstlisting}[style=csharp,basicstyle=\codexs,escapechar=|]
|\mk{f1}|services.AddScoped<IOrderRepository, SqlOrderRepository>();

|\mk{f2}|services.AddScoped<OrderService>();

|\mk{f3}|services.AddSingleton(TimeProvider.System);

|\mk{f4}|services.AddScoped<IOrderRepository>(provider =>
    new CachedOrderRepository(
        provider.GetRequiredService<SqlOrderRepository>()));

|\mk{f5}|services.TryAddSingleton<INotifier, EmailNotifier>();
\end{lstlisting}
  \begin{tikzpicture}[remember picture,overlay,every node/.style={rawlab}]
    \node[text=accent] at ([xshift=9.6cm]f1) {интерфейс и реализация};
    \node[text=accent] at ([xshift=9.6cm]f2) {только реализация};
    \node at ([xshift=9.6cm]f3) {готовый экземпляр};
    \node at ([xshift=9.6cm]f4) {фабрика};
    \node at ([xshift=9.6cm]f5) {если ещё не зарегистрировано};
  \end{tikzpicture}

  \slidegap
  \aside{Все формы про одно: как получить экземпляр. Почти всегда хватает
    \hi{первых двух}.}
\end{frame}

% 8
\begin{frame}{Регистрировать и разрешать}
  \slidetop
  \claim{Регистрация --- правило в списке. Разрешение --- \hi{дерево} по
    конструкторам.}

  \begin{center}
    \begin{tikzpicture}[font=\scriptsize,
      tb/.style={box,minimum width=2.25cm,minimum height=1.0cm,inner sep=2pt}]
      \node[tb,draw=accent,line width=1.0pt] (d1) at (0,0) {обработчик};
      \node[tb] (d2) at (2.85,0)  {\texttt{OrderService}};
      \node[tb] (d3) at (5.7,0)   {\texttt{IOrderRepository}\\\texttt{SqlOrderRepository}};
      \node[tb] (d4) at (8.55,0)  {\texttt{DbConnection}};
      \node[tb] (d5) at (11.4,0)  {\texttt{Options}};
      \foreach \i [evaluate=\i as \j using int(\i+1)] in {1,...,4}
        \draw[arr] (d\i) -- (d\j);
      \draw[arr] ([yshift=0.45cm]d1.north west) --
        node[above,text=muted] {строит сверху вниз} ([yshift=0.45cm]d5.north east);
      \draw[arra] ([yshift=-0.45cm]d5.south east) --
        node[below,text=accent] {создаёт снизу вверх} ([yshift=-0.45cm]d1.south west);
    \end{tikzpicture}
  \end{center}

  \slidegap
  \aside{План создания строится один раз и кэшируется: дерево из десяти
    объектов --- микросекунды. Разрешение вызывает фреймворк, вы --- почти
    никогда.}
\end{frame}

% 9
\begin{frame}{Конструктор не делает работу}
  \slidetop
  \versus{Конструктор}{%
    \begin{cbul}{accent}
      \item берёт зависимости
      \item кладёт их в поля
      \item создаётся в тесте без побочных эффектов
    \end{cbul}}{Не в конструкторе}{%
    \begin{cbul}{alt}
      \item открыть соединение
      \item прочитать файл, сходить в сеть
      \item запустить поток
    \end{cbul}}

  \slidegap
  \hrulethin{0.42}
  \rulegap
  \aside{Инициализация --- в явном асинхронном методе или в фоновом сервисе.
    Контейнер бесплатен, пока конструкторы \hi{только присваивают}.}
\end{frame}

% 10
\begin{frame}[plain]
  \statement{Конструктор --- единственная \hi{честная} точка входа
    зависимостей.}
\end{frame}

% 11
\begin{frame}[plain]
  \vote{Зачем интерфейс, если реализация одна и второй не планируется?}
\end{frame}

% 12
\begin{frame}{Времена жизни}
  \slidetop
  {\renewcommand{\arraystretch}{1.6}
  \begin{tabular}{@{}L{2.8cm}L{4.6cm}L{5.4cm}@{}}
    & \colhead{accent}{сколько живёт} & \colhead{alt}{пример} \\
    \textbf{Singleton} & один на приложение        & настройки, кэш, логгер \\
    \textbf{Scoped}    & один на запрос            & контекст базы, текущий
      пользователь \\
    \textbf{Transient} & новый на каждое разрешение & лёгкие объекты без
      состояния \\
  \end{tabular}}

  \slidegap
  \aside{Создаются \hi{лениво}, при первом обращении. Singleton --- один раз
    даже при одновременных первых обращениях из разных потоков.}
\end{frame}

% 13
\begin{frame}{Контейнер владеет тем, что создал}
  \slidetop
  \begin{center}
    \begin{tikzpicture}[font=\scriptsize]
      \node[box,draw=accent,line width=1.0pt,minimum width=12.4cm,minimum height=0.6cm,
            inner sep=1pt] (s) at (6.2,1.0) {singleton: от первого обращения до остановки приложения};
      \foreach \i/\x in {1/1.9,2/6.2,3/10.5} {
        \node[box,minimum width=3.4cm,minimum height=0.6cm,inner sep=1pt]
          (q\i) at (\x,0) {scoped: запрос \i};
        \node[text=alt,font=\scriptsize,anchor=north] at (q\i.south east)
          {\texttt{Dispose}};
      }
    \end{tikzpicture}
  \end{center}

  \slidegap
  \begin{cbul}{accent}
    \item scoped и transient освобождаются вместе с областью --- в конце запроса
    \item singleton --- при остановке приложения
    \item асинхронное освобождение контейнер вызывает сам
    \item пришло через конструктор --- \hi{вы не владелец}: \texttt{Dispose} не
      пишете
  \end{cbul}
\end{frame}

% 14
\begin{frame}[plain]
  \vote{Singleton-сервис принимает в конструктор scoped-зависимость. Что
    произойдёт?}
\end{frame}

% 15
\begin{frame}[fragile]{Captive dependency}
  \slidetop
  \begin{center}
    \begin{tikzpicture}[font=\footnotesize,
      tb/.style={box,minimum width=4.4cm,minimum height=1.1cm}]
      \node[tb] (c1) at (0,0)   {\textbf{singleton}\\живёт всегда};
      \node[tb,draw=alt,line width=1.0pt] (c2) at (7.6,0)
        {\textbf{scoped}\\должен жить один запрос};
      \draw[arrl] (c1) -- node[above,text=alt,font=\scriptsize\bfseries] {держит}
        node[below,text=muted,font=\scriptsize] {дольше, чем тот живёт} (c2);
    \end{tikzpicture}
  \end{center}

  \rulegap
\begin{lstlisting}[style=term,basicstyle=\termsm]
$ dotnet run 2>&1 | grep -o "Cannot consume[^.]*\." | head -1
Cannot consume scoped service 'IOrderRepository' from singleton 'PriceCache'.
\end{lstlisting}

  \rulegap
  \aside{Все следующие запросы работают с \lo{мёртвым} объектом первого.
    Проверка --- только в Development: в проде тот же код стартует молча.}
\end{frame}

% 16
\begin{frame}{Как выбрать время жизни}
  \slidetop
  \claim{Три вопроса к объекту.}

  {\renewcommand{\arraystretch}{1.6}
  \begin{tabular}{@{}L{5.6cm}L{7.2cm}@{}}
    есть изменяемое состояние?   & нет --- \hi{singleton}, самый дешёвый \\
    кому оно принадлежит?        & запросу --- \hi{scoped}; никому ---
      \hi{transient} \\
    держит ресурс?               & scoped; \lo{не} transient из корня \\
  \end{tabular}}

  \slidegap
  \aside{Хоть одна зависимость scoped --- сам объект тоже scoped. Контекст
    базы --- всегда scoped.}
\end{frame}

% 17
\begin{frame}{Пять аварий}
  \slidetop
  {\renewcommand{\arraystretch}{1.45}
  \begin{tabular}{@{}L{0.5cm}L{5.3cm}L{7.0cm}@{}}
    \lo{\textbf{1}} & singleton держит scoped        & один мёртвый объект на всё
      приложение \\
    \lo{\textbf{2}} & изменяемый singleton           & тысяча потоков пишут в один
      словарь \\
    \lo{\textbf{3}} & контекст базы не той длительности & transient --- два
      контекста; singleton --- аварии 1 и 2 \\
    \lo{\textbf{4}} & transient из корня             & не освобождается до
      остановки: утечка \\
    \lo{\textbf{5}} & \texttt{Dispose} чужого        & контейнер освободит второй
      раз \\
  \end{tabular}}

  \slidegap
  \aside{Что пришло через конструктор, освобождает контейнер. Что создали через
    \texttt{new} --- \hi{вы}.}
\end{frame}

% 18
\begin{frame}[fragile]{Несколько реализаций}
  \slidetop
\begin{lstlisting}[style=csharp,basicstyle=\codexs,escapechar=|]
services.AddScoped<IOrderCheck, StockCheck>();
services.AddScoped<IOrderCheck, LimitCheck>();

|\mk{m1}|public sealed class OrderChecks(IEnumerable<IOrderCheck> checks)

|\mk{m2}|public sealed class CachedOrderRepository(IOrderRepository inner)
    : IOrderRepository

|\mk{m3}|services.AddScoped(typeof(IRepository<>), typeof(SqlRepository<>));
\end{lstlisting}
  \begin{tikzpicture}[remember picture,overlay,every node/.style={rawlab}]
    \node[text=accent] at ([xshift=10.5cm]m1) {все, по порядку};
    \node at ([xshift=10.5cm]m2) {декоратор};
    \node at ([xshift=10.5cm]m3) {открытые обобщения};
  \end{tikzpicture}

  \slidegap
  \begin{cbul}{accent}
    \item «вызвать всех» --- в порядке регистрации: проверки, валидаторы
    \item декоратор --- обёртка снаружи: кэш, повтор; класс внутри \hi{не знает}
    \item одна регистрация на всё семейство хранилищ
  \end{cbul}
\end{frame}

% 19
\begin{frame}[fragile]{Сервисы с ключом}
  \slidetop
  \claim{«Вызвать \hi{конкретную}» --- ключ в сигнатуре.}

\begin{lstlisting}[style=csharp,basicstyle=\codexs,escapechar=|]
|\mk{k1}|builder.Services.AddKeyedScoped<INotifier, EmailNotifier>("email");
builder.Services.AddKeyedScoped<INotifier, SmsNotifier>("sms");

public sealed class OrderConfirmation(
|\mk{k2}|    [FromKeyedServices("email")] INotifier notifier)
{
    public Task SendAsync(Order order, CancellationToken ct) =>
        notifier.SendAsync(order.CustomerEmail,
            $"Order {order.Id} confirmed", ct);
}
\end{lstlisting}
  \begin{tikzpicture}[remember picture,overlay,every node/.style={rawlab}]
    \node at ([xshift=10.4cm]k1) {ключ при регистрации};
    \node[text=accent] at ([xshift=10.4cm]k2) {атрибут на параметре};
  \end{tikzpicture}

  \slidegap
  \aside{Ключ-перечисление безопаснее строки: опечатка --- ошибка компиляции.}
\end{frame}

% 20
\begin{frame}{Service locator}
  \slidetop
  \versus{Конструктор}{%
    \begin{cbul}{accent}
      \item видно всё, что нужно
      \item ошибка --- на старте
      \item тест подставляет заглушки
      \item времена жизни проверяет \texttt{Build}
    \end{cbul}}{IService Provider}{%
    \begin{cbul}{alt}
      \item зависимости спрятаны в теле
      \item ошибка --- на первом вызове
      \item тест собирает контейнер
      \item scoped из корня --- и никто не заметит
    \end{cbul}}

  \slidegap
  \hrulethin{0.42}
  \rulegap
  \aside{Законен \hi{на границе}: фабрика областей в фоновом сервисе, фильтр,
    middleware. В бизнес-логике --- сигнал на ревью.}
\end{frame}

% 21
\begin{frame}[plain]
  \vote{Кто писал «для удобства» статический класс, который внутри ходит в
    базу, читает файл или время?}
\end{frame}

% 22
\begin{frame}[plain]
  \vote{Сто сервисов --- сто строк с \texttt{Add} в файле запуска. Что делать?}
\end{frame}

% 23
\begin{frame}[fragile]{Модули}
  \vspace{0.3em}
\begin{lstlisting}[style=csharp,basicstyle=\codexxs,escapechar=|]
// Program.cs
|\mk{u1}|builder.Services
    .AddOrders()
    .AddNotifications()
    .AddPersistence(builder.Configuration);

// Orders/OrdersModule.cs
public static class OrdersModule
{
|\mk{u2}|    public static IServiceCollection AddOrders(
        this IServiceCollection services)
    {
|\mk{u3}|        services.AddScoped<IOrderRepository, SqlOrderRepository>();
        services.AddScoped<OrderService>();
        services.AddScoped<IValidator<CreateOrderDto>, CreateOrderDtoValidator>();
|\mk{u4}|        return services;
    }
}
\end{lstlisting}
  \begin{tikzpicture}[remember picture,overlay,every node/.style={rawlab}]
    \node[text=accent] at ([xshift=10.2cm]u1) {читается как оглавление};
    \node at ([xshift=10.2cm]u2) {метод расширения на модуль};
    \node at ([xshift=10.2cm]u3) {состав --- рядом с модулем};
    \node at ([xshift=10.2cm]u4) {чтобы шли цепочкой};
  \end{tikzpicture}
\end{frame}

% 24
\begin{frame}{Явно или по соглашению}
  \slidetop
  \versus{Явные модули}{%
    \begin{cbul}{accent}
      \item состав виден в одном месте
      \item модуль регистрирует только своё
      \item порядок вызовов не важен
      \item длиннее, зато ревьюится
    \end{cbul}}{Сканирование сборки}{%
    \begin{cbul}{alt}
      \item сотня строк исчезает
      \item состав --- по именам классов
      \item публичный класс регистрируется сам
      \item ограничивайте узко
    \end{cbul}}

  \slidegap
  \hrulethin{0.42}
  \rulegap
  \aside{Composition root стоит проверять \hi{тестом}: собрать контейнер и
    попросить у него каждый зарегистрированный тип.}
\end{frame}

% 25
\begin{frame}[plain]
  \keyline{Состав приложения --- в одном месте.\\
    Зависимости --- в конструкторе.\\
    Времена жизни --- контракт с контейнером.}
\end{frame}

% ==================================================== 2. OUTGOING CALLS ==
\section{2. Исходящие HTTP-вызовы}

% 26
\begin{frame}[plain]
  \sectionopener{2}{Исходящие HTTP-вызовы}{Внешний сервис --- это
    \hi{зависимость}, как хранилище.}
\end{frame}

% 27
\begin{frame}[plain]
  \vote{Нужно сделать HTTP-запрос из метода. Что вы пишете?}
\end{frame}

% 28
\begin{frame}{Две ошибки}
  \slidetop
  \versus{Клиент на каждый вызов}{%
    \begin{cbul}{accent}
      \item \texttt{new HttpClient} и \texttt{using}
      \item \texttt{using} закрывает соединение
      \item сотни соединений ждут закрытия
      \item сокеты \hi{кончились}
    \end{cbul}}{Один статический}{%
    \begin{cbul}{alt}
      \item создан один раз, навсегда
      \item соединение живёт, пока работает
      \item DNS разрешён один раз
      \item адрес сменился --- ходим на \lo{старый}
    \end{cbul}}

  \slidegap
  \hrulethin{0.42}
  \rulegap
  \aside{Клиент должен жить долго, чтобы переиспользовать соединения, но не
    вечно, чтобы они обновлялись.}
\end{frame}

% 29
\begin{frame}{Фабрика клиентов}
  \slidetop
  \begin{center}
    \begin{tikzpicture}[font=\scriptsize,
      tb/.style={box,minimum width=3.3cm,minimum height=1.0cm}]
      \node[tb] (h1) at (0,0)    {\textbf{клиенты}\\лёгкие, короткоживущие};
      \node[tb,draw=accent,line width=1.0pt] (h2) at (4.4,0)
        {\textbf{пул обработчиков}\\держат соединения};
      \node[tb] (h3) at (8.8,0)  {\textbf{внешний сервис}\\адрес из DNS};
      \draw[arr] (h1) -- (h2);
      \draw[arr] (h2) -- (h3);
      \node[text=accent,font=\scriptsize,anchor=north] at ([yshift=-0.15cm]h2.south)
        {заменяются раз в пару минут};
    \end{tikzpicture}
  \end{center}

  \slidegap
  \begin{cbul}{accent}
    \item клиент от фабрики создавайте сколько угодно
    \item обработчики делятся: соединения \hi{переиспользуются}
    \item обработчик заменён --- соединение переоткрыто, DNS перечитан
  \end{cbul}

  \rulegap
  \aside{Именованные клиенты --- строки вместо типов. Типизированные --- класс.}
\end{frame}

% 30
\begin{frame}[fragile]{Типизированный клиент}
  \vspace{0.3em}
\begin{lstlisting}[style=csharp,basicstyle=\codexxs,escapechar=`]
`\mk{r1}`public sealed class RatesClient(HttpClient http)
{
    public async Task<decimal> GetRateAsync(string currency, CancellationToken ct)
    {
        var response = await http.GetFromJsonAsync<RateResponse>($"rates/{currency}", ct);
        return response?.Rate
`\mk{r2}`            ?? throw new RatesUnavailableException(currency);
    }
}

builder.Services
`\mk{r3}`    .AddHttpClient<RatesClient>(client =>
    {
        client.BaseAddress = new Uri(builder.Configuration["Rates:BaseUrl"]
            ?? throw new InvalidOperationException("Rates:BaseUrl is not set"));
        client.Timeout = TimeSpan.FromSeconds(5);
    })
`\mk{r4}`    .AddStandardResilienceHandler();

`\mk{r5}`public sealed class PricingService(RatesClient rates) { /* ... */ }
\end{lstlisting}
  \begin{tikzpicture}[remember picture,overlay,every node/.style={rawlab}]
    \node[text=accent] at ([xshift=10.9cm]r1) {получает клиент};
    \node at ([xshift=10.9cm]r2) {своё исключение};
    \node at ([xshift=10.9cm]r3) {адрес и таймаут};
    \node at ([xshift=10.9cm]r4) {устойчивость};
    \node[text=accent] at ([xshift=10.9cm]r5) {обычная зависимость};
  \end{tikzpicture}
\end{frame}

% 31
\begin{frame}{Что живёт в клиенте}
  \slidetop
  \versus{Внутри клиента}{%
    \begin{cbul}{accent}
      \item адреса и параметры чужого API
      \item заголовки, которые он требует
      \item чужой DTO: \texttt{RateResponse}
      \item его ошибки --- ваши исключения
    \end{cbul}}{Снаружи}{%
    \begin{cbul}{alt}
      \item что делать с курсом
      \item потребитель не знает про HTTP
      \item ваш тип или примитив: \texttt{decimal}
      \item в тесте --- заглушка, сети нет
    \end{cbul}}

  \slidegap
  \hrulethin{0.42}
  \rulegap
  \aside{Три \texttt{new HttpClient} с адресом в строке превращаются в три
    параметра конструкторов.}
\end{frame}

% 32
\begin{frame}{Три инструмента}
  \slidetop
  \begin{center}
    \begin{tikzpicture}[font=\footnotesize,
      tb/.style={box,minimum width=4.0cm,minimum height=1.25cm}]
      \node[tb] (i1) at (0,0)    {\textbf{таймаут}\\не ждать вечно};
      \node[tb,draw=accent,line width=1.0pt] (i2) at (4.5,0)
        {\textbf{повтор}\\только идемпотентное};
      \node[tb] (i3) at (9.0,0)  {\textbf{выключатель}\\падать быстро,\\пока сосед лежит};
    \end{tikzpicture}
  \end{center}

  \slidegap
  \begin{cbul}{accent}
    \item таймаут --- секунды, а не умолчание в десятки секунд
    \item \texttt{GET} повторяется безопасно; \texttt{POST} --- \lo{только} с
      ключом идемпотентности
    \item выключатель считает ошибки и отвечает отказом сразу, без ожидания
  \end{cbul}

  \rulegap
  \aside{Повторы без паузы добивают перегруженного соседа: \lo{шторм повторов}.}
\end{frame}

% 33
\begin{frame}[plain]
  \vote{Внешний сервис ответил ошибкой. Сколько раз повторять?}
\end{frame}

% 34
\begin{frame}[fragile]{Стандартный набор}
  \slidetop
\begin{lstlisting}[style=csharp,basicstyle=\codesm]
builder.Services.AddHttpClient<RatesClient>(/* ... */)
    .AddStandardResilienceHandler();
\end{lstlisting}

  \slidegap
  \begin{center}
    \begin{tikzpicture}[font=\scriptsize,
      tb/.style={box,minimum width=2.7cm,minimum height=0.9cm,inner sep=2pt}]
      \node[tb] (t1) at (0,0)    {общий таймаут};
      \node[tb] (t2) at (3.2,0)  {повтор};
      \node[tb] (t3) at (6.4,0)  {выключатель};
      \node[tb] (t4) at (9.6,0)  {таймаут попытки};
      \foreach \i [evaluate=\i as \j using int(\i+1)] in {1,...,3}
        \draw[arr] (t\i) -- (t\j);
    \end{tikzpicture}
  \end{center}

  \slidegap
  \begin{cbul}{accent}
    \item таймаутов \hi{два}: на одну попытку и на всю операцию с повторами
    \item попыток немного, пауза растёт, разброс случайный
    \item начинайте с умолчаний: они \hi{согласованы} между собой
  \end{cbul}
\end{frame}

% 35
\begin{frame}[fragile]{Когда сосед лежит}
  \slidetop
  {\renewcommand{\arraystretch}{1.5}
  \begin{tabular}{@{}L{3.4cm}L{9.4cm}@{}}
    \hi{\textbf{курс валюты}} & взять из кэша и пометить приблизительным: сервис
      деградировал, но работает \\
    \lo{\textbf{платёж}}      & честная ошибка «попробуйте позже»: 503 \\
  \end{tabular}}

  \slidegap
\begin{lstlisting}[style=term,basicstyle=\termsm]
info: System.Net.Http.HttpClient.RatesClient.LogicalHandler[100]
      Start processing HTTP request GET http://localhost:5090/stub/rates/USD
...
info: System.Net.Http.HttpClient.RatesClient.LogicalHandler[101]
      End processing HTTP request after 103.371ms - 200
\end{lstlisting}

  \rulegap
  \aside{Решение принимает ваш сервис, и \hi{заранее}. Каждый исходящий вызов
    фабрика пишет в лог: длительность и статус.}
\end{frame}

% 36
\begin{frame}[plain]
  \keyline{\texttt{HttpClient} выдаёт фабрика; внешний сервис приходит в
    конструктор.\\ Повторять можно только то, что безопасно повторить.}
\end{frame}

% ================================================== 3. CONFIGURATION ==
\section{3. Конфигурация и окружение}

% 37
\begin{frame}[plain]
  \sectionopener{3}{Конфигурация и окружение}{Код один, окружений много;
    разница живёт \hi{снаружи}.}
\end{frame}

% 38
\begin{frame}[plain]
  \vote{Какую ошибку в первом листинге лекции мы исправили не до конца?}
\end{frame}

% 39
\begin{frame}{Что считать конфигурацией}
  \slidetop
  \versus{Конфигурация}{%
    \begin{cbul}{accent}
      \item адреса, строки подключения, порты
      \item лимиты и таймауты
      \item уровни логов
      \item переключатели возможностей
    \end{cbul}}{Код}{%
    \begin{cbul}{alt}
      \item бизнес-правила
      \item как считается скидка
      \item выбор класса по флагу
      \item всё, что нужно тестировать
    \end{cbul}}

  \slidegap
  \hrulethin{0.42}
  \rulegap
  \aside{Граница: значение может поменять \hi{администратор}, и от этого ничего
    не сломается.}
\end{frame}

% 40
\begin{frame}{Слои}
  \slidetop
  \begin{center}
    \begin{tikzpicture}[font=\scriptsize,
      lay/.style={box,minimum width=5.6cm,minimum height=0.56cm,inner sep=1pt},
      txt/.style={anchor=west,text=muted,font=\scriptsize,inner sep=0pt}]
      \node[lay] (l1) at (0,0)    {\texttt{appsettings.json}};
      \node[lay] (l2) at (0,0.68) {\texttt{appsettings.\{Environment\}.json}};
      \node[lay] (l3) at (0,1.36) {секреты пользователя};
      \node[lay] (l4) at (0,2.04) {переменные окружения};
      \node[lay] (l5) at (0,2.72) {аргументы командной строки};
      \node[lay,draw=accent,line width=1.0pt] (l6) at (0,3.6) {хранилище секретов};
      \node[txt] at ([xshift=0.4cm]l1.east) {умолчания};
      \node[txt] at ([xshift=0.4cm]l2.east) {перекрывает по своим ключам};
      \node[txt] at ([xshift=0.4cm]l3.east) {только в разработке, вне репозитория};
      \node[txt] at ([xshift=0.4cm]l4.east) {окружение процесса};
      \node[txt] at ([xshift=0.4cm]l5.east) {\texttt{args}};
      \node[txt,text=accent] at ([xshift=0.4cm]l6.east) {в проде, подключается последним};
      \draw[arra] ([xshift=-0.5cm]l1.south west) --
        node[left,text=accent,align=right,font=\scriptsize] {позднее\\перекрывает\\раннее}
        ([xshift=-0.5cm]l6.north west);
    \end{tikzpicture}
  \end{center}

  \rulegap
  \aside{Каждый источник --- плоский набор пар «ключ --- значение». Побеждает
    верхний, у которого ключ есть.}
\end{frame}

% 41
\begin{frame}[fragile]{Три записи одного ключа}
  \slidetop
\begin{lstlisting}[style=json,basicstyle=\codexs]
{
  "ConnectionStrings": { "Orders": "Host=localhost;Database=orders" },
  "Rates": { "BaseUrl": "https://rates.local/", "TimeoutSeconds": 5 },
  "Logging": { "LogLevel": { "Default": "Information" } }
}
\end{lstlisting}

\begin{lstlisting}[style=bash,basicstyle=\codexs]
export ConnectionStrings__Orders="Host=db.prod.internal;Database=orders"
export Rates__TimeoutSeconds=2
dotnet HelloWorld.dll --Logging:LogLevel:Default=Warning
\end{lstlisting}

  \rulegap
  {\renewcommand{\arraystretch}{1.25}\small
  \begin{tabular}{@{}L{3.6cm}L{9.2cm}@{}}
    в приложении          & \hi{\texttt{Rates:TimeoutSeconds}} \\
    JSON                  & вложенные объекты \\
    переменная окружения  & \texttt{Rates\_\_TimeoutSeconds} --- два подчёркивания \\
    аргумент              & \texttt{-{}-Rates:TimeoutSeconds=2} \\
  \end{tabular}}
\end{frame}

% 42
\begin{frame}[fragile]{Что видит приложение}
  \slidetop
  \claim{Отладочное представление: ключ, значение и \hi{откуда} оно пришло.}

\begin{lstlisting}[style=term,basicstyle=\termsm]
$ curl -s localhost:5090/config | grep -E "Orders=|Default=|BaseUrl=|TimeoutSeconds="
  Orders=Host=db.prod.internal;Database=orders (EnvironmentVariablesConfigurationProvider)
    Default=Warning (CommandLineConfigurationProvider)
  BaseUrl=https://rates.local/ (JsonConfigurationProvider for 'appsettings.json' (Optional))
  TimeoutSeconds=2 (EnvironmentVariablesConfigurationProvider)
\end{lstlisting}

  \slidegap
  \begin{cbul}{accent}
    \item нет ключа --- \texttt{null}, а \lo{не} исключение: опечатка молчит
    \item по ключу читает файл запуска; бизнес-логика --- \lo{нет}
    \item в этом выводе будут и секреты: в прод --- никогда
  \end{cbul}
\end{frame}

% 43
\begin{frame}{Окружения}
  \slidetop
  \claim{Имя окружения --- строка в \hi{переменной}.}

  {\renewcommand{\arraystretch}{1.6}
  \begin{tabular}{@{}L{4.4cm}L{8.4cm}@{}}
    \hi{\textbf{второй файл}}          & \texttt{appsettings.Production.json}
      читается только в проде \\
    \hi{\textbf{поведение фреймворка}} & страница ошибок, проверка контейнера,
      сертификат разработки \\
    \lo{\textbf{ваши развилки}}        & «если разработка, то» --- чем меньше, тем
      лучше \\
  \end{tabular}}

  \slidegap
  \aside{Локально имя задаёт файл настроек запуска. В опубликованное приложение
    он не попадает.}
\end{frame}

% 44
\begin{frame}{Разница --- в значениях}
  \slidetop
  \versus{В конфигурации}{%
    \begin{cbul}{accent}
      \item адрес шлюза --- из конфигурации
      \item в разработке там адрес заглушки
      \item код один и тот же
      \item в проде работает то, что тестировали
    \end{cbul}}{В коде}{%
    \begin{cbul}{alt}
      \item «если прод --- настоящий шлюз»
      \item «иначе --- заглушка»
      \item две ветки
      \item в проде работает другая ветка
    \end{cbul}}

  \slidegap
  \hrulethin{0.42}
  \rulegap
  \aside{Стенд похож на прод во всём, кроме адресов. Проверка окружения в
    домене --- тот же запах, что \texttt{HttpContext} в домене.}
\end{frame}

% 45
\begin{frame}[plain]
  \vote{Переменная с именем окружения не задана вовсе. Каким приложение себя
    считает?}
\end{frame}

% 46
\begin{frame}{Рантайм настраивается так же}
  \slidetop
  {\renewcommand{\arraystretch}{1.22}\small
  \begin{tabular}{@{}L{6.2cm}L{6.6cm}@{}}
    \texttt{ASPNETCORE\_ENVIRONMENT}          & окружение \\
    \texttt{ASPNETCORE\_URLS}                 & адреса Kestrel \\
    \texttt{Logging\_\_LogLevel\_\_Default}   & уровень логов \\
    \hi{\texttt{DOTNET\_gcServer}}            & серверный GC (Л2) \\
    \hi{\texttt{DOTNET\_GCHeapHardLimit}}     & потолок памяти (Л2) \\
    \hi{\texttt{DOTNET\_ThreadPool\_MinThreads}} & минимум потоков пула (Л6) \\
    \hi{\texttt{DOTNET\_TieredCompilation}}   & многоуровневая компиляция (Л1) \\
    \hi{\texttt{DOTNET\_ROLL\_FORWARD}}       & накат версии рантайма (Л2) \\
  \end{tabular}}

  \slidegap
  \aside{\hi{\texttt{DOTNET\_}} читает рантайм, до первой строки вашего кода.
    Сервис в проде ведёт себя иначе --- сравните \hi{окружение}, а не код.}
\end{frame}

% 47
\begin{frame}[plain]
  \vote{Один образ нужно запустить на стенде и в проде, с разными базами. Что
    делать с файлом настроек?}
\end{frame}

% 48
\begin{frame}{Образ один}
  \slidetop
  \begin{center}
    \begin{tikzpicture}[font=\footnotesize,
      env/.style={box,minimum width=2.6cm,minimum height=0.8cm},
      var/.style={font=\codexxs,text=muted,fill=paper,inner sep=2pt}]
      \node[box,draw=accent,line width=1.0pt,minimum width=2.4cm,minimum height=1.4cm]
        (img) at (0,0) {\textbf{образ}\\собран один раз};
      \node[env] (e1) at (10.2,1.35)  {разработка};
      \node[env] (e2) at (10.2,0)     {стенд};
      \node[env] (e3) at (10.2,-1.35) {прод};
      \draw[arr] (img.east) -- node[var,pos=0.55] {Host=localhost} (e1.west);
      \draw[arr] (img.east) -- node[var,pos=0.55] {Host=db.stage.internal} (e2.west);
      \draw[arr] (img.east) -- node[var,pos=0.55] {Host=db.prod.internal} (e3.west);
      \node[font=\scriptsize,text=muted,anchor=north] at (0,-0.85)
        {\texttt{ConnectionStrings\_\_Orders}};
    \end{tikzpicture}
  \end{center}

  \slidegap
  {\large Образ один, среды разные, разница --- в \hi{переменных}.}

  \rulegap
  \aside{Ровно то, что протестировали, едет в прод. Секретов в образе нет: ни в
    файле, ни в коде, ни в истории.}
\end{frame}

% 49
\begin{frame}{Почему переменные}
  \slidetop
  \versus{Переменные окружения}{%
    \begin{cbul}{accent}
      \item процесс просто их видит
      \item задаёт всё: оркестратор, скрипт, CI
      \item существуют только у процесса
    \end{cbul}}{Файл в образе}{%
    \begin{cbul}{alt}
      \item монтирование, пути, права
      \item его нужно куда-то положить
      \item живёт в слоях образа навсегда
    \end{cbul}}

  \slidegap
  \hrulethin{0.42}
  \rulegap
  \aside{Файл остаётся: умолчания и структура. Монтируют большие списки и
    секреты. Переменные \lo{видно} в списке процессов и описании развёртывания.}
\end{frame}

% 50
\begin{frame}[plain]
  \vote{Где хранить строку подключения к локальной базе, чтобы она не попала в
    репозиторий?}
\end{frame}

% 51
\begin{frame}[fragile]{Секреты пользователя}
  \slidetop
\begin{lstlisting}[style=term,basicstyle=\termsm]
$ dotnet user-secrets set "ConnectionStrings:Orders" "Host=localhost;Password=dev"
Successfully saved ConnectionStrings:Orders to the secret store.
$ curl -s localhost:5090/config | grep -A1 "^ConnectionStrings"
ConnectionStrings:
  Orders=Host=localhost;Password=dev (JsonConfigurationProvider for 'secrets.json' (Optional))
\end{lstlisting}

  \slidegap
  \begin{cbul}{accent}
    \item \texttt{dotnet user-secrets init} --- файл в домашней папке, привязан к проекту
    \item читается \hi{только} в разработке, в репозитории не лежит
    \item в \texttt{appsettings.json} --- структура и безобидные умолчания
  \end{cbul}

  \rulegap
  \aside{Проверка: репозиторий завтра станет публичным --- что утечёт?}
\end{frame}

% 52
\begin{frame}[plain]
  \keyline{Разница между окружениями живёт в конфигурации.\\
    Образ один, сред много: сверху --- переменные.}
\end{frame}

% ==================================================== 4. OPTIONS ==
\section{4. Options pattern}

% 53
\begin{frame}[plain]
  \sectionopener{4}{Options pattern}{Настройки --- это \hi{тип}, а не строка
    по ключу.}
\end{frame}

% 54
\begin{frame}[fragile]{Класс настроек}
  \vspace{0.3em}
\begin{lstlisting}[style=csharp,basicstyle=\codexxs,escapechar=|]
public sealed class RatesOptions
{
|\mk{p1}|    public const string Section = "Rates";

|\mk{p2}|    [Required, Url]
    public string BaseUrl { get; set; } = "";

    [Range(1, 60)]
|\mk{p3}|    public int TimeoutSeconds { get; set; } = 5;
}

builder.Services
    .AddOptions<RatesOptions>()
|\mk{p4}|    .BindConfiguration(RatesOptions.Section)
|\mk{p5}|    .ValidateDataAnnotations()
|\mk{p6}|    .ValidateOnStart();

public sealed class PricingService(
|\mk{p7}|    IOptionsMonitor<RatesOptions> options) { /* ... */ }
\end{lstlisting}
  \begin{tikzpicture}[remember picture,overlay,every node/.style={rawlab}]
    \node at ([xshift=10.6cm]p1) {имя секции --- в классе};
    \node at ([xshift=10.6cm]p2) {правила атрибутами};
    \node at ([xshift=10.6cm]p3) {единица --- в имени};
    \node at ([xshift=10.6cm]p4) {привязать к секции};
    \node at ([xshift=10.6cm]p5) {проверить атрибуты};
    \node[text=accent] at ([xshift=10.6cm]p6) {упасть на старте};
    \node at ([xshift=10.6cm]p7) {обычная зависимость};
  \end{tikzpicture}
\end{frame}

% 55
\begin{frame}{Три интерфейса}
  \slidetop
  \claim{Разница в одном: \hi{когда} значение перечитывается.}

  {\renewcommand{\arraystretch}{1.6}
  \begin{tabular}{@{}L{4.2cm}L{5.0cm}L{3.6cm}@{}}
    \hi{\texttt{IOptions<T>}}     & один раз, до перезапуска        & singleton \\
    \texttt{IOptionsSnapshot<T>}  & один раз на запрос              & scoped: в
      singleton \lo{нельзя} \\
    \texttt{IOptionsMonitor<T>}   & всегда актуальное, с уведомлением & singleton \\
  \end{tabular}}

  \slidegap
  \aside{По умолчанию --- \texttt{IOptions}. Перечитываются файлы; переменные
    окружения процесс читает один раз, при старте.}
\end{frame}

% 56
\begin{frame}[fragile]{Клиент из настроек}
  \slidetop
  \claim{Адрес живёт в \hi{одном} классе. Строковых ключей нет.}

\begin{lstlisting}[style=csharp,basicstyle=\codexs,escapechar=|]
builder.Services
|\mk{w1}|    .AddHttpClient<RatesClient>((provider, client) =>
    {
        RatesOptions options = provider
|\mk{w2}|            .GetRequiredService<IOptions<RatesOptions>>().Value;
        client.BaseAddress = new Uri(options.BaseUrl);
        client.Timeout = TimeSpan.FromSeconds(options.TimeoutSeconds);
    })
    .AddStandardResilienceHandler();
\end{lstlisting}
  \begin{tikzpicture}[remember picture,overlay,every node/.style={rawlab}]
    \node at ([xshift=10.4cm]w1) {перегрузка с контейнером};
    \node[text=accent] at ([xshift=10.4cm]w2) {проверенные настройки};
  \end{tikzpicture}

  \slidegap
  \aside{Контейнер в регистрации --- это граница, а не бизнес-логика.}
\end{frame}

% 57
\begin{frame}[plain]
  \vote{Вы забыли задать адрес курсов в проде. Когда вы об этом узнаете?}
\end{frame}

% 58
\begin{frame}[fragile]{Падать на старте}
  \slidetop
  \versus{С Validate On Start}{%
    \begin{cbul}{accent}
      \item процесс не поднялся
      \item трафик остался на старой версии
      \item ошибка в логе запуска, в ту же минуту
    \end{cbul}}{Без него}{%
    \begin{cbul}{alt}
      \item стартует, отвечает на \texttt{/health}
      \item принимает трафик
      \item падает первый запрос за курсом, через час
    \end{cbul}}

  \slidegap
\begin{lstlisting}[style=term,basicstyle=\termsm]
$ Rates__BaseUrl= dotnet run 2>&1 | grep -m1 -o "for 'RatesOptions.*"
for 'RatesOptions' members: 'BaseUrl' with the error: 'The BaseUrl field is required.'.
\end{lstlisting}

  \rulegap
  \aside{Сообщение пишите так, чтобы было понятно, какую переменную задать.}
\end{frame}

% 59
\begin{frame}[plain]
  \keyline{Настройки --- тип, а не строка по ключу.\\
    Проверяются на старте, пока ошибка ничего не стоит.}
\end{frame}

\end{document}
